% !TEX program = xelatex
%
% korean-test.tex — Smoke test for latex-korean-setup skill.
%
% Build with:
%   xelatex korean-test.tex
%
% Expected: PDF renders all Korean glyphs correctly. If you see boxes (□)
% or missing characters, your TeX Live installation is missing kotex or
% the engine is misconfigured.
%
\documentclass[11pt]{article}

% kotex은 TeX Live 기본 패키지이므로 추가 설치 불필요.
% Overleaf, MacTeX, Linux TeX Live 모두에서 동일하게 작동.
\usepackage[hangul]{kotex}

\title{한글 LaTeX 렌더링 검증}
\author{latex-korean-setup skill}
\date{\today}

\begin{document}
\maketitle

\section{서론}
이 문서는 LaTeX 환경이 한글을 정상적으로 렌더링할 수 있는지 확인하기 위한
smoke test 파일입니다. 이 PDF가 모든 한글 글자를 깨짐 없이 표시한다면
환경 셋업이 성공한 것입니다.

\section{검증 항목}

\subsection{기본 한글 본문}
가나다라마바사아자차카타파하. 한글의 모든 자모가 정상적으로 표시되어야 합니다.
받침이 있는 글자 (\textbf{값}, \textbf{닭}, \textbf{흙}) 와 복모음 (\textbf{왜},
\textbf{외}, \textbf{위}) 도 깨지지 않아야 합니다.

\subsection{한영 혼용}
한글과 English가 섞인 문장은 LaTeX (라텍)으로 작성된 academic paper
(학술 논문) 에서 매우 흔합니다. 두 글자체 사이 간격이 자연스러워야 합니다.

\subsection{기호와 숫자}
숫자: 0, 1, 2, \ldots, 100. 기호: $\alpha, \beta, \gamma, \pi$. 한글 기호:
※, ★, ☆, ◯, ●. 따옴표: ``큰따옴표'', `작은따옴표'.

\subsection{목록}
\begin{itemize}
  \item 첫 번째 항목 — em dash 처리
  \item 두 번째 항목 (괄호 포함)
  \item 세 번째 항목: 콜론 뒤 한글
\end{itemize}

\begin{enumerate}
  \item 가
  \item 나
  \item 다
\end{enumerate}

\subsection{표}
\begin{tabular}{|l|c|r|}
\hline
이름 & 점수 & 비고 \\
\hline
가갸 & 100 & 우수 \\
나냐 & 85  & 양호 \\
다댜 & 72  & 보통 \\
\hline
\end{tabular}

\section{수식과 한글}
한글 본문 중간에 수식 $E = mc^2$ 이 들어가도 정렬이 맞아야 합니다.
displayed equation도 마찬가지:
\[
  \int_0^\infty e^{-x^2}\,dx = \frac{\sqrt{\pi}}{2}
\]
수식 다음에도 한글이 자연스럽게 이어져야 합니다.

\section{결론}
이 페이지의 모든 글자가 \textbf{읽기 가능한 형태로} 표시되었다면, 한글 LaTeX
환경 셋업이 정상 완료된 것입니다. 만약 빈 박스(\texttt{□}) 나 깨진 글자가
보인다면 \texttt{kotex} 패키지 또는 xelatex 엔진이 누락된 것입니다.

문제 진단:
\begin{verbatim}
python ~/.claude/skills/latex-korean-setup/scripts/diagnose.py
\end{verbatim}

\end{document}
